\documentclass[11pt,a4paper]{article}
\usepackage[utf8]{inputenc}
\usepackage[spanish]{babel}
\usepackage{amsmath}
\usepackage{amsfonts}
\usepackage{amssymb}
\usepackage[margin=2.5cm]{geometry}
\usepackage{xcolor}
\usepackage{graphicx} % MUY IMPORTANTE: Necesario para que funcione \resizebox

\begin{document}

\begin{center}
    \Large \textbf{Etapa de Estimación y Predicción} \\[0.5cm]
    \normalsize \textbf{Estadística Aplicada Resumen de Estudio} \\
    \textbf{Roberto Francisco Espinola} \\
    \textbf{Septiembre 2026}
\end{center}

\vspace{0.5cm}

Habiendo comprobado que nuestro modelo es significativo, es decir, al rechazar la hipótesis nula de que $\beta = 0$, y al haber relación lineal entre $X$ e $Y$, podemos estimar valores de la variable dependiente $Y$. A continuación, analizaremos qué ocurre si queremos realizar estimaciones fijando la variable independiente en un punto específico, por ejemplo, $x_0 = 190 \text{ cm/seg}$. O sea, vamos a estimar el valor promedio de todas las turbinas que operen a $190 \text{ cm/s}$.

\vspace{0.4cm}

\subsection*{Paso 4: Intervalo de Confianza para la Media de Y ($x_0 = 190$)}

El estimador puntual es:
\[
\hat{y}_0 = a + b \cdot 190 = 0.069242 + 0.003829 \cdot 190 = \mathbf{0.7967} \text{ mm}^2\text{/seg}
\]

El error máximo de estimación para la media poblacional $\mu_y$ con un 95\% de confianza es:
\[
E = t_{\alpha/2, n-2} \cdot s_e \cdot \sqrt{\frac{1}{n} + \frac{n(x_0 - \bar{x})^2}{S_{xx}}}
\]

Calculamos la amplitud del margen de error sustituyendo los datos en la fórmula de la raíz:
\[
E = 2.306 \cdot 0.159052 \cdot \sqrt{\frac{1}{10} + \frac{10(190 - 200)^2}{1,320,000}}
\]
\[
E = 0.36677 \cdot \sqrt{0.1 + \frac{10(-10)^2}{1,320,000}}
\]
\[
E = 0.36677 \cdot \sqrt{0.1 + \frac{1000}{1,320,000}} = 0.36677 \cdot \sqrt{0.100757} \approx \textbf{0.1164}
\]

\vspace{0.3cm}
Sustituyendo los valores numéricos consolidados, la inecuación queda configurada de la siguiente manera:

\[
\resizebox{\textwidth}{!}{$\displaystyle
\underbrace{0.7967}_{\hat{y}_0} - \underbrace{\left( 2.306 \cdot 0.159052 \cdot \sqrt{\frac{1}{10} + \frac{10(190 - 200)^2}{1,320,000}} \right)}_{E \approx 0.1164} < \mu_y < \underbrace{0.7967}_{\hat{y}_0} + \underbrace{\left( 2.306 \cdot 0.159052 \cdot \sqrt{\frac{1}{10} + \frac{10(190 - 200)^2}{1,320,000}} \right)}_{E \approx 0.1164}
$}
\]

\vspace{0.3cm}
Resolviendo la resta en el lado izquierdo y la suma en el lado derecho, obtenemos los límites:
\begin{itemize}
    \item \textbf{Extremo Inferior:} $0.7967 - 0.1164 = 0.6803$
    \item \textbf{Extremo Superior:} $0.7967 + 0.1164 = 0.9131$
\end{itemize}

\vspace{0.3cm}
\noindent \textbf{Conclusión del Intervalo:}
\[
\mathbf{0.6803 < \mu_y < 0.9131}
\]

\vspace{0.8cm}
\hrule
\vspace{0.8cm}

Y ahora, ¿qué pasa si queremos predecir qué pasará si $x_0 = 190$, o sea, predecir una turbina en particular a una velocidad de 190?

\vspace{0.6cm}
\subsection*{Paso 5: Intervalo de Predicción para una Observación Futura Individual ($x_0 = 190$)}

¿Qué pasa si en lugar de predecir la media de muchos ensayos queremos anticipar el resultado de \textbf{un único ensayo futuro}? La incertidumbre será mucho mayor, puesto que a la variabilidad de estimar la recta se le suma la dispersión física natural del proceso ($\sigma^2$). La fórmula del error máximo de predicción individual agrega un ``1'' dentro de la raíz:

\[
\resizebox{\textwidth}{!}{$\displaystyle
\underbrace{\hat{y}_0}_{\text{Estimador puntual}} - \underbrace{t_{\alpha/2, n-2} \cdot s_e \cdot \sqrt{1 + \frac{1}{n} + \frac{n(x_0 - \bar{x})^2}{S_{xx}}}}_{E_{pred} \text{ (Margen de Error)}} < y_0 < \underbrace{\hat{y}_0}_{\text{Estimador puntual}} + \underbrace{t_{\alpha/2, n-2} \cdot s_e \cdot \sqrt{1 + \frac{1}{n} + \frac{n(x_0 - \bar{x})^2}{S_{xx}}}}_{E_{pred} \text{ (Margen de Error)}}
$}
\]

\vspace{0.2cm}
\noindent Calculamos la amplitud del margen de error sustituyendo los datos en la fórmula de la raíz:

\[
E_{pred} = 2.306 \cdot 0.159052 \cdot \sqrt{1 + \frac{1}{10} + \frac{10(190 - 200)^2}{1,320,000}}
\]
\[
E_{pred} = 0.36677 \cdot \sqrt{1 + 0.1 + \frac{10(-10)^2}{1,320,000}}
\]
\[
E_{pred} = 0.36677 \cdot \sqrt{1 + 0.1 + \frac{1000}{1,320,000}} = 0.36677 \cdot \sqrt{1.100757} \approx \mathbf{0.3848}
\]

\vspace{0.4cm}

\[
\resizebox{\textwidth}{!}{$\displaystyle
\underbrace{0.7967}_{\hat{y}_0} - \underbrace{\left( 2.306 \cdot 0.159052 \cdot \sqrt{1 + \frac{1}{10} + \frac{10(190 - 200)^2}{1,320,000}} \right)}_{E_{pred} \approx 0.3848} < y_0 < \underbrace{0.7967}_{\hat{y}_0} + \underbrace{\left( 2.306 \cdot 0.159052 \cdot \sqrt{1 + \frac{1}{10} + \frac{10(190 - 200)^2}{1,320,000}} \right)}_{E_{pred} \approx 0.3848}
$}
\]

\vspace{0.4cm}
\noindent Resolviendo la resta en el lado izquierdo y la suma en el lado derecho, obtenemos los límites:

\begin{itemize}
    \item \textbf{Extremo Inferior:} $0.7967 - 0.3848 = 0.4119$
    \item \textbf{Extremo Superior:} $0.7967 + 0.3848 = 1.1815$
\end{itemize}

\vspace{0.3cm}
\noindent \textbf{Conclusión del Intervalo:}
\[
\mathbf{0.4119 < y_0 < 1.1815}
\]

Podemos predecir, con un 95\% de confianza, que si extraemos una turbina individual y fijamos la velocidad del aire a $190 \text{ cm/seg}$, su coeficiente de evaporación específico se encontrará entre $0.4119 \text{ y } 1.1815 \text{ mm}^2\text{/seg}$.

\vspace{0.5cm}
\noindent \textcolor{blue}{\textbf{Nota de análisis estadístico (Tip para el final oral):}} \\
Si el tribunal te pide justificar por qué este intervalo ($[0.4119, 1.1815]$) es más ancho que el de la media ($[0.6803, 0.9131]$), la respuesta técnica es que predecir un evento individual incluye \textbf{dos fuentes de error}. No solo estamos lidiando con la incertidumbre de dónde está posicionada exactamente la recta verdadera (error de los estimadores $a$ y $b$), sino que también debemos sumarle la variabilidad inherente del sistema (la varianza poblacional $\sigma^2$), porque un ensayo individual puede caer en cualquier parte de la campana de Gauss alrededor de esa recta. Ese "1" extra dentro de la raíz representa matemáticamente esa dispersión física.

\end{document}
